{
\AMCquestionNumberfalse
\def\AMCbeginQuestion#1#2{}
\noindent

\begin{question}{Questão 1}
Qual é o valor aproximado de \pi?

\begin{multicols}{4}\columnseprule=0.1pt
\begin{choices}
\correctchoice{3.14}
\wrongchoice{4.13}
\wrongchoice{6.28}
\wrongchoice{1.41}
\end{choices}
\end{multicols}
\end{question}
}

{
\AMCquestionNumberfalse
\def\AMCbeginQuestion#1#2{}
\noindent

\begin{question}{Questão 2}
Qual das opções abaixo contém o conjunto da intersecção de {1, 2, 3, 5, 7} e {1, 3, 9}

\begin{multicols}{4}\columnseprule=0.1pt
\begin{choices}
\correctchoice{1, 3}
\wrongchoice{2, 4}
\wrongchoice{5, 7}
\wrongchoice{2, 9}
\end{choices}
\end{multicols}
\end{question}
}

{
\AMCquestionNumberfalse
\def\AMCbeginQuestion#1#2{}
\noindent

\begin{question}{Questão 3}
Considere um triângulo retângulo com lados iguais a 6 cm e 8 cm. Assinale a(s) alternativa(s) correta(s):

\begin{multicols}{4}\columnseprule=0.1pt
\begin{choices}
\correctchoice{A hipotenusa vale 10 cm}
\correctchoice{A área do triângulo é 24 cm^2}
\wrongchoice{É um triângulo equilátero}
\wrongchoice{O perímetro é 22 cm}
\wrongchoice{A área do triângulo é 48 cm^2}
\end{choices}
\end{multicols}
\end{question}
}